\documentclass[11pt,a4paper]{article}
\usepackage[utf8]{inputenc}
\usepackage[T1]{fontenc}
\usepackage{estilo_capstone}

\titulo{Product Vision — Evento Cultural}
\subtitulo{Visión de producto (artefacto 1 del marco ágil, anexo metodológico)}
\autor{Roberto Riffo · José Miguel Moncada · David Herrera}
\fecha{25 de septiembre de 2026}

\newcolumntype{L}[1]{>{\raggedright\arraybackslash}p{#1}}

\begin{document}
\maketitlecapstone

\section{Problema}

Los eventos culturales, comerciales y municipales de Chile se publican de forma
dispersa: redes sociales (principalmente Instagram), páginas municipales y medios
locales. No existe un catálogo único, normalizado y siempre vigente donde una
persona pueda descubrir qué hacer; la información llega tarde, incompleta o
desaparece con el evento. Este problema está documentado en el design brief del
proyecto y respaldado por las historias de usuario HU-01 a HU-15
(\code{Historias\_de\_Usuario\_Plataforma\_Cultural.xlsx}), derivadas de problemas
reales expresados por personas en fuentes públicas.

\section{Usuarios}

\begin{table}[h!]
\centering\small
\begin{tabular}{@{}L{4.2cm}L{10.6cm}@{}}
\toprule
\textbf{Usuario} & \textbf{Necesidad principal (evidencia)} \\
\midrule
Visitante & Consultar una cartelera web sin depender de redes sociales (HU-01),
conocer eventos antes de que terminen (HU-02), planificar por fecha (HU-03),
conocer precio y gratuidad (HU-04) y ubicación (HU-05). \\
Usuario registrado & Recibir recomendaciones según sus intereses (HU-15) y
reportar errores del catálogo (HU-13). \\
Organizador & Publicar su evento en el catálogo (flujo de organización presente
en el frontend: página \code{/publicar}). \\
\bottomrule
\end{tabular}
\caption{Usuarios y necesidades respaldadas por el backlog de historias de usuario.}
\end{table}

\section{Propuesta de valor}

Evento Cultural es una plataforma que compila y difunde eventos de todo el país
en un catálogo único, normalizado y siempre vigente:

\begin{itemize}
  \item \textbf{Recolección automática:} un scraper en Python integra agendas
        públicas desde 21 fuentes, con normalización, filtro cultural,
        deduplicación y checkpoints reanudables; la ejecución se automatiza con
        GitHub Actions.
  \item \textbf{Catálogo canónico:} Supabase (PostgreSQL) centraliza los eventos
        con bitácora de ejecuciones y políticas de Row Level Security.
  \item \textbf{Descubrimiento:} frontend Angular 22 + Ionic 9 (web y móvil) con
        búsqueda, filtros, mapa, calendario y recomendaciones por señales de
        interés.
  \item \textbf{Origen visible:} cada evento declara su fuente y su última
        comprobación, para que la persona confirme antes de asistir (HU-12).
\end{itemize}

\section{Qué hace diferente a la solución}

\begin{itemize}
  \item \textbf{Sin colecciones permanentes:} no existen favoritos ni guardados;
        los eventos rotan, vencen y desaparecen. Las acciones del usuario se
        expresan como señales de interés temporales que alimentan las
        recomendaciones.
  \item \textbf{Foco cultural y local:} filtro cultural en el pipeline y
        cobertura de actividades municipales y de región, que hoy no tienen
        difusión centralizada (HU-14).
  \item \textbf{Trazabilidad del dato:} origen y vigencia declarados en cada
        registro, a diferencia de las carteleras informales en redes sociales.
\end{itemize}

\section{Alcance de la versión actual}

El estado real del sistema a la fecha: scraper operativo con ejecución manual y
automatizada en GitHub Actions (evidencia: captura de ejecución del workflow),
base de datos en Supabase con esquema y RLS documentados, y frontend con
catálogo mock en memoria (la conexión real a Supabase es un pendiente
declarado). El sistema no usa Docker: se aloja y automatiza con GitHub Actions
(scraper, desplegado y en testeo del despliegue), Vercel (hosting de la página,
en preparación) y Capacitor 8 + Android Studio (la web se convertirá en
aplicación Android).

\begin{pendienteenv}
métricas de éxito del producto (adquisición, retención, eventos publicados por
organizadores): deben definirlas y medirlas el equipo.
\end{pendienteenv}

\end{document}
